% !TEX program = xelatex
% AI-effects 프로젝트 --- AI 점유율 지표 A~C(paper 27)와 노출도 지표 3개의 분포 비교(분산, 쏠림, 분포 간 거리, 탄력성, 중심까지의 거리).
% 표 본문: results/table/33_disp.tex, 33_hhi.tex, 33_conc.tex, 33_curve.tex, 33_bdiff.tex, 33_dist.tex, 33_ksco1.tex,
%          33_elas.tex, 33_dsum.tex, 33_d_pos.tex, 33_d_neg.tex, 33_top_C.tex (code/04_16_analysis_kor_share_exposure_dist.do가 생성)
% 전체 자료: results/table/33_kor_share_exposure_dist.xlsx
% 관련: paper 25(RU1·RU2), paper 27(점유율 지표), paper 28(순위상관·노출도 기반 점유율), paper 29(순위 격차·로그 비율·회귀 잔차), paper 32(규모 탄력성)
% 컴파일: xelatex -> biber -> xelatex -> xelatex (kotex, biblatex/biber 필요)
\documentclass[11pt]{article}

\usepackage{fontspec}
\usepackage{kotex}
\usepackage[a4paper,margin=2.0cm]{geometry}
\usepackage{longtable}
\usepackage{booktabs}
\usepackage{array}
\usepackage{amsmath}
\usepackage{xcolor}
\usepackage{hyperref}
\usepackage{enumitem}
\usepackage[backend=biber,style=authoryear,maxcitenames=2,uniquelist=false,sorting=nyt]{biblatex}
\addbibresource{../../reference/references.bib}

\hypersetup{colorlinks=true,linkcolor=blue!50!black,urlcolor=blue!50!black}
\newcolumntype{L}[1]{>{\raggedright\arraybackslash}p{#1}}
\newcolumntype{R}[1]{>{\raggedleft\arraybackslash}p{#1}}
\setlength{\LTleft}{0pt}
\setlength{\LTright}{0pt}
\renewcommand{\arraystretch}{1.15}
\setlength{\tabcolsep}{3.5pt}

\title{AI 점유율 지표와 노출도 지표의 분포 비교\\[2mm]
\large --- 순위가 담지 못하는 분산, 쏠림, 분포 간 거리 ---}
\author{AI-effects 프로젝트}
\date{\today}

\begin{document}
\maketitle

\begin{abstract}
\noindent paper 28$\cdot$29는 점유율 지표 A$\sim$C와 노출도 지표(Observed exposure, Felten AIOE, Eloundou GPT-4 $\beta$)를 주로 순위로 비교했다. 순위는 값이 얼마나 퍼져 있는지, 어디에 몰려 있는지를 버린다. 이 노트는 두 지표를 같은 척도에 놓고 분포 자체를 비교한다. 척도는 둘이다. 하나는 취업자 가중 평균이 1인 상대 강도(사용 $R$, 노출도 $RX$)이고, 다른 하나는 합이 1인 점유율(사용 $s$, 노출도 기반 점유율 $XS$)이다. 그 위에서 네 가지를 보았다. 분산(지니, Theil), 노출도를 따라 사용이 쏠리는 정도(집중지수, 비유사성 지수, KL 거리), 탄력성, 중심까지의 거리다. 표본은 KSCO 8차 소분류 163개다(2026-11-01부터 한국 대화 비중이 두 달 모두 미공개인 34개를 사용 0으로 넣는다). 로그를 쓰는 값(로그 분산, KL 거리, 탄력성, 로그 거리, 로버스트 $z$)은 사용 $> 0$인 129개로 계산한다. 구간은 소분류를 복원추출한 부트스트랩 999회의 95\% 백분위 구간이다.
\begin{itemize}[leftmargin=*]
  \item \textbf{사용은 노출도보다 훨씬 고르지 않다.} 상대 강도의 취업자 가중 지니는 사용 0.87$\sim$0.90, Observed exposure 0.64, Eloundou $\beta$ 0.37이다. 차이(+0.24$\sim$+0.53)의 구간은 모두 0을 포함하지 않는다. 퍼짐의 차이는 위쪽보다 아래쪽 꼬리에서 크다. 90분위는 사용 1.6$\sim$1.7배, $\beta$ 2.0배, OE 3.0배로 큰 차이가 없다. 하지만 사용은 소분류 79개(사용 0인 34개 포함)가 평균의 10분의 1에도 못 미치고, $\beta$로는 2개뿐이다.
  \item \textbf{사용의 직업 간 차이는 대분류 안에서 더 많이 생긴다.} Theil 지수 중 대분류 간 몫은 사용 40$\sim$43\%, 노출도 68$\sim$74\%다. 노출도는 ``사무$\cdot$전문가 대 현장직''의 차이가 대부분이다. 사용은 같은 대분류 안에서도 직업마다 크게 다르다.
  \item \textbf{노출도를 따라 쏠리는 정도는 노출도 지표에 따라 다르다.} 노출도 순으로 줄 세운 사용의 집중지수는 0.58$\sim$0.67이다. Observed exposure로 만든 노출도 기반 점유율의 집중지수(0.64)와는 차이가 없다(+0.02$\sim$+0.03, 구간이 0 포함). Eloundou $\beta$(0.37)보다는 0.27$\sim$0.29 크다(구간 0.11$\sim$0.45). 탄력성도 같은 그림이다. 사용 $> 0$인 129개에서 $\ln R$를 $\ln RX$에 회귀한 기울기는 Observed exposure에 대해 1.00$\sim$1.07이다(사용 강도가 노출도에 비례). Eloundou $\beta$에 대해서는 1.80$\sim$2.21이다($\beta$의 차이가 사용에서 두 배 넘게 증폭).
  \item \textbf{집중의 크기는 같아도 위치가 다르다.} 점유율 분포의 실효 직업 수(1/HHI)는 지표 C 22.0, 노출도 점유율(OE) 24.6으로 거의 같다. 그런데 사용 분포를 노출도 분포와 같게 만들려면 사용 점유율의 60$\sim$69\%를 다른 직업으로 옮겨야 한다(비유사성 지수). 그 차이의 약 70\%는 대분류 사이의 이동이다. 전문가 대분류는 사용 점유율이 71.9\%인데 노출도 기반 점유율은 29$\sim$35\%다. 사무$\cdot$판매는 노출도 기반 점유율이 44$\sim$56\%인데 사용 점유율은 16\%다.
  \item \textbf{거리로 보면 순위와 다른 직업이 드러난다.} 지표 C와 Eloundou $\beta$의 로그 거리 $d$는 사용 $> 0$인 129개 중 108개에서 두 배를 넘고, 51개에서 열 배를 넘는다. $d$와 순위 격차의 순위상관은 0.26뿐이다. $d$는 취업자 수와 음의 상관($-0.57$)을 보인다. 큰 직업일수록 사용이 노출도에 못 미친다는 뜻이다. 다만 $d$의 아래쪽 끝은 사용 점유율이 거의 0인 직업(조리사, 작물 재배, 자동차 운전원 등)이 차지한다. 이 값들은 추정이 가장 불안정한 곳이므로 로버스트 $z$ 차이 $\Delta z$와 함께 읽어야 한다. 사용이 0인 34개는 $d$와 $\Delta z$가 정의되지 않는다.
\end{itemize}
\end{abstract}

\section{비교 방법}

\subsection{같은 척도}

점유율 지표(paper 27)는 합이 100인 규모 지표이고, 노출도는 0$\sim$1 값(Observed exposure, Eloundou $\beta$)이거나 평균 0 근처의 표준화 점수(Felten)인 강도 지표다. 둘의 값을 바로 빼거나 평균까지의 거리를 비교하면 단위가 맞지 않는다. 그래서 먼저 둘을 같은 척도로 옮긴다. 이 노트에서 쓰는 모든 비중은 소분류 163개 안에서 다시 정규화한 값이다(사용 0인 34개는 사용 점유율 0).

소분류 $j$의 취업자 수를 $E_j$(2022년 상반기 지역별고용조사), 취업자 비중을 $w_j = E_j / \sum_k E_k$라 하자. 지표 $m \in \{\text{A}, \text{B}, \text{C}\}$의 점유율을 $s^m_j$(합 1)라 하고, 노출도 $x \in \{\text{OE}, \beta\}$의 노출도 기반 점유율(paper 28)을 다음과 같이 둔다.
\[
XS^x_j = \frac{E_j x_j}{\sum_k E_k x_k}.
\]
그러면 두 가지 같은 척도가 생긴다.
\begin{itemize}[leftmargin=*]
  \item \textbf{점유율(규모)}: $s^m_j$와 $XS^x_j$. 둘 다 합이 1이다.
  \item \textbf{상대 강도}: 사용 $R^m_j = s^m_j / w_j$와 노출도 $RX^x_j = XS^x_j / w_j = x_j / \bar{x}_E$다. $\bar{x}_E$는 노출도의 취업자 가중 평균이다. 둘 다 취업자 가중 평균이 1이므로 ``평균의 몇 배인가''로 읽는다. $R^{\mathrm{A}}$는 RU1에, $R^{\mathrm{C}}$는 RU2에 비례한다(paper 25).
\end{itemize}
$s^m_j = R^m_j w_j$, $XS^x_j = RX^x_j w_j$이므로 두 척도의 규모 부분은 같고, 차이는 강도 부분에서만 생긴다. Felten AIOE는 음수가 있어 $XS$와 $RX$를 만들 수 없다. 그래서 Felten은 순서만 쓰는 방법(집중지수)과 로버스트 $z$에만 넣었다. 사용 0인 34개는 $\ln R$이, Observed exposure가 0인 28개(그중 13개는 사용이 있음)는 $\ln RX$가 정의되지 않는다. 그래서 로그를 쓰는 통계는 사용 $> 0$인 129개(OE와 짝지으면 116개)로 계산한다.

\subsection{방법 2: 분산}

상대 강도 $R^m$, $RX^x$ 각각에 대해 퍼짐을 수치 하나로 요약한다. 모두 취업자 가중이다.
\begin{itemize}[leftmargin=*]
  \item 변동계수 CV(가중 평균이 1이므로 가중 표준편차와 같음), 로그 분산 $\mathrm{var}(\ln R)$, 분위수 p10$\cdot$p90과 그 비.
  \item 지니계수: $R$ 오름차순으로 누적 취업자 비중 대비 누적 점유율을 그린 로렌츠 곡선으로 계산한다.
  \item Theil T: $T = \sum_j s_j \ln(s_j / w_j)$. KSCO 대분류 $g$로 $T = \sum_g S_g \ln(S_g / W_g) + \sum_g S_g T_g$와 같이 나눌 수 있다. 앞 항(대분류 간)의 몫으로 차이가 어느 수준에서 생기는지 본다.
\end{itemize}
노출도는 강도가 퍼진 정도가 원래 다를 수 있다. 예를 들어 $\beta$는 값이 0.01$\sim$0.92에 고르게 퍼져 있다. 따라서 ``사용이 노출도보다 X배 더 퍼져 있다''는 노출도 지표마다 다르게 나온다. 이 점은 순위 비교에서는 보이지 않던 차이다.

\subsection{방법 3: 쏠림}

\begin{itemize}[leftmargin=*]
  \item \textbf{집중곡선과 집중지수}: 소분류를 노출도 $x$ 오름차순으로 줄 세운다. 가로축에 누적 취업자 비중, 세로축에 누적 사용 점유율 $s^m$을 그린 것이 집중곡선이다. 집중지수는 $\mathrm{CI} = 1 - 2 \times (\text{곡선 아래 면적})$로 정의한다. 사용이 취업자에 비례하면 CI는 0이고, 노출도가 높은 직업에 몰릴수록 1에 가깝다. 같은 노출도 값(OE = 0인 28개 등)은 한 점으로 묶었다. 비교 기준은 노출도 기반 점유율의 집중지수 $\mathrm{CI}(XS^x)$다. 이 값은 정의상 $RX^x$의 지니계수와 같다(do 파일에서 assert). $\mathrm{CI}(s^m) > \mathrm{CI}(XS^x)$이면 ``노출도가 예측하는 것보다 사용이 고노출 직업에 더 몰려 있다''는 뜻이다. 또 $\mathrm{CI}(s^m) / \mathrm{Gini}(R^m)$은 사용의 불평등 중 노출도 순서와 같은 방향인 몫이다. 1이면 사용 순서가 노출도 순서와 같다.
  \item \textbf{노출도 상위 $q$의 몫}: 노출도가 가장 높은 취업자 $q$(10, 20, 50\%)가 차지하는 사용 점유율과 노출도 기반 점유율이다. 집중곡선의 몇 점을 읽은 것이다.
  \item \textbf{비유사성 지수}: $\mathrm{DI} = \tfrac{1}{2} \sum_j |s^m_j - XS^x_j|$. ``사용 분포를 노출도 분포와 같게 만들려면 사용 점유율의 몇 \%를 다른 직업으로 옮겨야 하는가''다. 대분류 합으로만 계산한 DI를 대분류 간 몫으로 함께 보인다(대분류 간 DI $\le$ 전체 DI).
  \item \textbf{KL 거리}: $D(XS^x \,\|\, s^m) = \sum_j XS^x_j \ln(XS^x_j / s^m_j)$. Theil 지수와 같은 형태라 대분류 간$\cdot$대분류 내로 정확히 나뉜다. 방향을 노출도 기준으로 둔 것은 OE = 0인 소분류 때문이다. 반대 방향 $D(s \,\|\, XS^{\mathrm{OE}})$는 무한대가 되고, $D(XS \,\|\, s)$는 사용 $> 0$인 129개 안에서 $XS$를 다시 정규화해 계산하므로 유한하다. $\beta$에 대한 반대 방향 값은 엑셀에 있다.
  \item \textbf{집중도}: 점유율 분포 자체의 HHI $\sum_j s_j^2$, 실효 직업 수 $1/\mathrm{HHI}$, 상위 10개 합.
\end{itemize}
DI와 KL은 무게를 두는 곳이 다르다. DI는 점유율 차이의 절댓값을 더하므로 몫이 큰 직업 사이의 이동에 민감하다. KL은 로그 비율에 무게를 두므로, 노출도 기반 점유율은 있는데 사용이 거의 없는 직업에 민감하다.

\subsection{방법 4: 탄력성}

\[
\ln R^m_j = \alpha + \beta^{m,x} \ln RX^x_j + \varepsilon_j
\]
를 취업자 가중 OLS(robust SE)로 추정한다. $\beta^{m,x} = 1$이면 사용 강도가 노출도에 비례한다. 1보다 크면 노출도의 차이가 사용에서 증폭되고(쏠림), 1보다 작으면 완화된다. 비가중 OLS와 분위수 회귀($\tau$ = 0.1, 0.5, 0.9)도 함께 보았다. paper 29의 방법 4는 $\ln s$를 노출도 수준값과 $\ln E$에 회귀한 잔차로 직업을 찾는 것이었다. 여기서는 양변을 상대 강도의 로그로 바꿔 기울기 자체를 ``증폭률''로 읽는다. Felten은 $\ln R$를 원값에 회귀한 준탄력성만 엑셀에 둔다.

\subsection{방법 1: 중심까지의 거리}

\begin{itemize}[leftmargin=*]
  \item \textbf{로그 거리}: $d^{m,x}_j = \ln R^m_j - \ln RX^x_j$. 두 지표 모두 중심(가중 평균)이 1이므로, 사용이 평균의 몇 배인지와 노출도가 평균의 몇 배인지를 비교한 것이다. $e^{d}$는 ``사용 강도가 노출도 강도의 몇 배인가''다. 정의상 $d = \ln(s^m / XS^x)$여서 paper 29의 로그 비율 LR와 같은 값이다(assert). 이 노트에서는 이 값을 직업 목록을 만드는 주 지표로 쓴다.
  \item \textbf{로버스트 $z$ 차이}: $z(v) = (v - \mathrm{med}(v)) / (1.4826\,\mathrm{MAD}(v))$, $\Delta z^{m,x}_j = z(\ln R^m_j) - z(\text{노출도}_j)$다. 노출도는 OE$\cdot\beta$이면 $\ln RX$, Felten이면 원값이고, 비가중이다. 중위값과 MAD로 각자의 퍼짐을 맞춘 뒤 비교하므로, $d$와 달리 두 분포의 퍼짐 차이를 지운다. 즉 $d$는 ``얼마나 떨어져 있나''(크기)를, $\Delta z$는 ``각자의 분포 안에서 위치가 얼마나 다른가''를 본다. Felten도 비교할 수 있다.
\end{itemize}

\subsection{불확실성}

소분류 163개를 복원추출해 같은 계산을 999번 반복했다(seed 20261025). 매번 취업자 비중, 점유율, 노출도 기반 점유율을 표본 안에서 다시 정규화한다. 표의 대괄호는 95\% 백분위 구간이다. 소분류 안의 표본 오차(지역별고용조사, AEI 비중)는 반영하지 않는다. 따라서 이 구간은 ``어떤 소분류가 표본에 들었느냐''에 대한 민감도로 읽어야 한다.

\section{결과}

\subsection{분산: 사용은 노출도보다 훨씬 고르지 않다}

\begin{small}
\begin{longtable}{L{4.3cm}R{1.0cm}R{1.2cm}R{1.0cm}R{1.0cm}R{1.4cm}R{1.9cm}R{1.9cm}R{1.3cm}}
\caption{상대 강도의 분산(취업자 가중, 소분류 163개; 로그 분산은 사용 $> 0$인 129개). 대괄호는 부트스트랩 95\% 구간. p10이 0이라 p90/p10은 사용과 OE에서 정의되지 않는다. $^{\ast}$ OE = 0인 13개를 더 뺀 116개}\label{tab:disp}\\
\toprule
지표 & CV & 로그 분산 & p10 & p90 & p90/p10 & 지니 & Theil T & 대분류 간(\%) \\
\midrule
\endfirsthead
% 04_16_analysis_kor_share_exposure_dist.do가 생성. 직접 고치지 말 것.
\(R^{\mathrm{A}}\) 대화 & 3.68 & 7.51 & 0.001 & 2.18 & 2424.2 & 0.85 & 1.85 & 37 \\
 & & & & & & {\footnotesize [0.79, 0.89]} & {\footnotesize [1.48, 2.21]} & \\
\(R^{\mathrm{B}}\) 절감 시간 & 3.91 & 10.59 & 0.000 & 2.05 & 27348.2 & 0.89 & 2.02 & 40 \\
 & & & & & & {\footnotesize [0.84, 0.92]} & {\footnotesize [1.63, 2.43]} & \\
\(R^{\mathrm{C}}\) 취업자 환산 & 3.90 & 10.61 & 0.000 & 1.99 & 23551.4 & 0.88 & 2.01 & 39 \\
 & & & & & & {\footnotesize [0.84, 0.91]} & {\footnotesize [1.61, 2.42]} & \\
\midrule
\(RX^{\mathrm{OE}}\) Observed exposure & 1.15 & 2.21$^{\ast}$ & 0.000 & 2.68 & -- & 0.60 & 0.63 & 65 \\
 & & & & & & {\footnotesize [0.49, 0.69]} & {\footnotesize [0.43, 0.87]} & \\
\(RX^{\beta}\) Eloundou \(\beta\) & 0.62 & 0.69 & 0.358 & 2.05 & 5.7 & 0.34 & 0.19 & 73 \\
 & & & & & & {\footnotesize [0.28, 0.39]} & {\footnotesize [0.13, 0.25]} & \\

\bottomrule
\end{longtable}

\begin{longtable}{L{3.0cm}L{1.0cm}R{1.6cm}R{2.4cm}R{1.6cm}R{2.4cm}}
\caption{사용과 노출도의 지니$\cdot$집중지수 차이(사용 $-$ 노출도)와 부트스트랩 95\% 구간}\label{tab:bdiff}\\
\toprule
사용 지표 & 노출도 & 지니 차이 & 95\% 구간 & 집중지수 차이 & 95\% 구간 \\
\midrule
\endfirsthead
% 04_16_analysis_kor_share_exposure_dist.do가 생성. 직접 고치지 말 것.
A 대화 & OE & +0.24 & [0.16, 0.31] & +0.02 & [-0.10, 0.12] \\
 & $\beta$ & +0.50 & [0.44, 0.56] & +0.27 & [0.11, 0.41] \\
\midrule
B 절감 시간 & OE & +0.26 & [0.19, 0.35] & +0.03 & [-0.08, 0.13] \\
 & $\beta$ & +0.53 & [0.48, 0.59] & +0.29 & [0.13, 0.45] \\
\midrule
C 취업자 환산 & OE & +0.26 & [0.19, 0.35] & +0.03 & [-0.08, 0.13] \\
 & $\beta$ & +0.53 & [0.47, 0.59] & +0.29 & [0.12, 0.44] \\

\bottomrule
\end{longtable}
\end{small}

\paragraph{퍼짐의 크기.} 사용 강도의 지니는 0.87$\sim$0.90, Theil은 1.99$\sim$2.16이다. 노출도는 Observed exposure가 지니 0.64, Theil 0.73이고, Eloundou $\beta$가 0.37, 0.22다(표~\ref{tab:disp}). 사용과 노출도의 지니 차이는 OE에 대해 +0.24$\sim$+0.26, $\beta$에 대해 +0.50$\sim$+0.53이고, 구간이 모두 0에서 멀다(표~\ref{tab:bdiff}). 사용 0인 34개가 들어가 사용의 지니가 129개만 보던 이전 판(0.85$\sim$0.89)보다 조금 커졌다. 범위로 보면 차이가 더 분명하다. 사용이 있는 소분류의 $R^{\mathrm{C}}$는 0.00001(451 조리사)에서 86(213 생명 및 자연과학 관련 시험원)까지 퍼져 있고, 34개는 0이다. 반면 $RX^{\beta}$는 0.017$\sim$2.88, $RX^{\mathrm{OE}}$는 0$\sim$5.9다. 순위 비교에서는 이 차이가 보이지 않는다. 순위상관 0.73$\sim$0.75(paper 28)는 ``순서가 꽤 비슷하다''는 뜻일 뿐, 사용의 직업 간 격차가 노출도의 격차보다 훨씬 크다는 사실은 담지 못한다.

\paragraph{어느 꼬리인가.} p90은 사용 1.62$\sim$1.74, $\beta$ 1.99, OE 2.95로 큰 차이가 없다. 차이는 아래쪽 꼬리에 있다. 사용의 p10은 0이고(사용 0인 34개가 취업자의 12.6\%), $R^{\mathrm{C}} < 0.1$인 소분류가 79개, $< 0.01$인 소분류가 46개(사용 0인 34개 포함)다. $\beta$로는 $RX < 0.1$인 소분류가 2개뿐이고 p90/p10은 6.6배다. 노출도 $\beta$는 ``현장직도 과업의 일부는 LLM에 노출된다''고 보지만, 실제 업무용 대화는 그런 직업에서 거의 나오지 않는다. OE는 소분류 47개가 $RX < 0.1$이고 28개가 0이라 아래 꼬리가 사용 쪽에 더 가깝다. 실제 사용 자료를 일부 쓴 지표이기 때문이다\parencite{massenkoff-2026-labor-market-impacts-of-ai}.

\paragraph{어느 수준에서 생기는가.} Theil 지수의 대분류 간 몫은 사용 40$\sim$43\%, OE 68\%, $\beta$ 74\%다. 노출도의 직업 간 차이는 대부분 ``사무$\cdot$전문가 대 현장직''이라는 대분류의 차이다\parencite{eloundou-2023-gpts-are-gpts-an-early}. 사용은 차이의 절반 넘게가 같은 대분류 안에서 생긴다. 예를 들어 전문가 대분류 안에서도 213 생명 및 자연과학 관련 시험원($R^{\mathrm{C}}$ = 86)과 244 간호사(0.011)의 차이가 약 8천 배다. 노출도 지표는 대분류 안의 이런 차이를 거의 만들지 못한다.

\subsection{쏠림: 노출도를 따라 얼마나 몰려 있는가}

\begin{small}
\begin{longtable}{L{3.8cm}R{2.1cm}R{2.1cm}R{2.1cm}R{2.4cm}R{2.4cm}}
\caption{노출도 순으로 줄 세운 집중지수. 대괄호는 부트스트랩 95\% 구간}\label{tab:conc}\\
\toprule
정렬 노출도 & CI(A) & CI(B) & CI(C) & CI($XS^x$) = 지니($RX^x$) & CI(C) / 지니($R^{\mathrm{C}}$) \\
\midrule
\endfirsthead
% 04_16_analysis_kor_share_exposure_dist.do가 생성. 직접 고치지 말 것.
Observed exposure & 0.66 & 0.67 & 0.67 & 0.64 & 0.75 \\
  & {\footnotesize [0.51, 0.76]} & {\footnotesize [0.52, 0.77]} & {\footnotesize [0.53, 0.77]} & {\footnotesize [0.54, 0.71]} & \\
Felten LM AIOE & 0.58 & 0.60 & 0.61 & -- & 0.67 \\
  & {\footnotesize [0.40, 0.71]} & {\footnotesize [0.43, 0.74]} & {\footnotesize [0.43, 0.74]} & & \\
Eloundou GPT-4 $\beta$ & 0.64 & 0.66 & 0.66 & 0.37 & 0.73 \\
  & {\footnotesize [0.48, 0.77]} & {\footnotesize [0.50, 0.79]} & {\footnotesize [0.49, 0.78]} & {\footnotesize [0.31, 0.41]} & \\

\bottomrule
\end{longtable}

\begin{longtable}{L{3.8cm}R{2.2cm}R{1.8cm}R{1.8cm}R{1.8cm}R{2.6cm}}
\caption{노출도 상위 $q$ 취업자가 차지하는 몫(\%)}\label{tab:curve}\\
\toprule
정렬 노출도 & 상위 $q$(\%) & A & B & C & $XS^x$ \\
\midrule
\endfirsthead
% 04_16_analysis_kor_share_exposure_dist.do가 생성. 직접 고치지 말 것.
Observed exposure & 10 & 51.6 & 52.3 & 51.9 & 38.6 \\
 & 20 & 67.1 & 67.3 & 67.2 & 63.4 \\
 & 50 & 94.4 & 95.4 & 95.4 & 95.4 \\
\midrule
Felten LM AIOE & 10 & 34.3 & 35.8 & 36.5 & -- \\
 & 20 & 41.9 & 44.3 & 44.8 & -- \\
 & 50 & 92.7 & 93.4 & 93.3 & -- \\
\midrule
Eloundou GPT-4 $\beta$ & 10 & 43.8 & 44.0 & 43.4 & 24.6 \\
 & 20 & 70.2 & 69.3 & 69.0 & 42.1 \\
 & 50 & 94.4 & 96.2 & 96.0 & 76.4 \\

\bottomrule
\end{longtable}

\begin{longtable}{L{6.0cm}R{2.0cm}R{2.6cm}R{2.6cm}}
\caption{점유율 분포의 집중도(소분류 163개)}\label{tab:hhi}\\
\toprule
분포 & HHI & 실효 직업 수 & 상위 10개 합(\%) \\
\midrule
\endfirsthead
% 04_16_analysis_kor_share_exposure_dist.do가 생성. 직접 고치지 말 것.
지표 A 대화 점유율 & 0.045 & 22.4 & 54.7 \\
지표 B 절감 시간 점유율 & 0.045 & 22.1 & 55.5 \\
지표 C 취업자 환산 점유율 & 0.045 & 22.0 & 55.5 \\
\midrule
노출도 점유율(OE) & 0.041 & 24.6 & 56.1 \\
노출도 점유율($\beta$) & 0.026 & 39.0 & 41.3 \\
\midrule
취업자 비중 & 0.019 & 53.8 & 33.5 \\

\bottomrule
\end{longtable}
\end{small}

\paragraph{집중지수.} 세 노출도 어느 순서로 줄 세워도 사용은 고노출 직업에 강하게 몰려 있다(CI 0.58$\sim$0.67, 표~\ref{tab:conc}). 비교 기준과 견주면 결론이 노출도에 따라 갈린다.
\begin{itemize}[leftmargin=*]
  \item \textbf{Observed exposure}: 노출도 기반 점유율의 집중지수가 0.64다. 사용의 집중지수(0.66$\sim$0.67)와 차이가 없다(+0.02$\sim$+0.03, 구간 $-0.10\sim+0.13$). 노출도 순서에 따른 쏠림의 크기는 OE가 예측하는 만큼이다.
  \item \textbf{Eloundou $\beta$}: 기준이 0.37이라 사용의 쏠림(0.64$\sim$0.66)이 0.27$\sim$0.29 더 크다(구간 0.11$\sim$0.45). $\beta$가 예측하는 것보다 실제 사용은 고노출 쪽에 훨씬 더 몰려 있다.
  \item \textbf{Felten}: 기준을 만들 수 없지만, CI가 0.58$\sim$0.61로 셋 중 가장 낮다.
\end{itemize}
CI(C) / 지니($R^{\mathrm{C}}$)는 0.67$\sim$0.75다. 사용의 불평등 중 약 3분의 2 이상이 노출도 순서와 같은 방향이고, 나머지는 노출도 순서와 어긋나는 재배열이다.

\paragraph{상위 몫.} 표~\ref{tab:curve}는 쏠림이 어디서 생기는지 보여 준다. OE 상위 10\% 취업자는 사용 점유율의 52\%를 차지한다. 노출도 기반 점유율로는 39\%다. 즉 OE의 맨 위에서 사용이 OE의 예측보다 더 몰린다. 대신 OE 하위 절반에서는 사용이 5$\sim$6\%만 남아 노출도의 예측(5\%)과 비슷하다. $\beta$는 상위 10\%에서 사용 43$\sim$44\%, 노출도 기반 점유율 25\%다. 상위 20\%에서는 69$\sim$70\% 대 42\%, 상위 50\%에서는 94$\sim$96\% 대 76\%로 벌어진다. $\beta$ 하위 절반의 취업자는 $\beta$ 기준으로 노출의 24\%를 갖지만 사용은 4$\sim$6\%뿐이다. 129개만 보던 이전 판에서는 $\beta$ 상위 10\%의 사용 몫이 23\%로 노출도 기반 점유율과 같았다. $\beta$ 최상위에 사용이 적은 큰 사무직(311 기획 및 마케팅 사무원, 331 회계 및 경리 사무원)이 있기 때문이었다. 취업자 분모가 163개로 커지면서 상위 10\% 구간이 넓어져, 사용이 많은 291 작가 및 언론 관련 전문가($\beta$ 10위, 지표 C 11.1\%)가 그 안에 들어왔다. 이 비교는 구간 경계에 민감하다.

\paragraph{집중의 크기와 위치.} 점유율 분포 자체의 집중도는 사용과 OE가 거의 같다(표~\ref{tab:hhi}). 실효 직업 수는 지표 C 22.0, 노출도 점유율(OE) 24.6이고, 상위 10개 합은 55.5\% 대 56.1\%다. $\beta$ 기반 점유율은 39.0개로 덜 몰려 있고, 취업자 분포(53.8개)에 가깝다. 즉 ``얼마나 몰려 있나''만 보면 사용과 OE는 같은 모양이다. 아래의 분포 간 거리는 둘이 서로 다른 직업에 몰려 있음을 보여 준다.

\subsection{분포 간 거리: 몰린 곳이 다르다}

\begin{small}
\begin{longtable}{L{2.4cm}L{0.8cm}R{3.0cm}R{1.3cm}R{2.8cm}R{1.3cm}R{1.7cm}R{1.7cm}}
\caption{비유사성 지수(DI)와 KL 거리 $D(XS \,\|\, s)$. 대괄호는 부트스트랩 95\% 구간, ``대분류 간''은 대분류 합으로 계산한 값의 몫. 오른쪽 두 열은 취업자 분포와의 DI}\label{tab:dist}\\
\toprule
사용 지표 & 노출도 & DI(\%) & 대분류 간(\%) & KL & 대분류 간(\%) & DI(사용, 취업자) & DI($XS$, 취업자) \\
\midrule
\endfirsthead
% 04_16_analysis_kor_share_exposure_dist.do가 생성. 직접 고치지 말 것.
A 대화 & OE & 59.9 {\footnotesize [47.1, 68.8]} & 67 & 1.35 {\footnotesize [0.96, 1.68]} & 31 & 72.7 & 51.1 \\
 & $\beta$ & 64.4 {\footnotesize [55.6, 71.5]} & 72 & 1.81 {\footnotesize [1.28, 2.33]} & 31 & 72.7 & 27.3 \\
\midrule
B 절감 시간 & OE & 63.4 {\footnotesize [49.4, 72.9]} & 68 & 1.73 {\footnotesize [1.22, 2.17]} & 35 & 77.1 & 51.1 \\
 & $\beta$ & 69.3 {\footnotesize [60.4, 76.2]} & 72 & 2.31 {\footnotesize [1.69, 2.91]} & 34 & 77.1 & 27.3 \\
\midrule
C 취업자 환산 & OE & 62.9 {\footnotesize [48.7, 73.0]} & 68 & 1.74 {\footnotesize [1.22, 2.18]} & 34 & 77.1 & 51.1 \\
 & $\beta$ & 68.8 {\footnotesize [59.9, 76.2]} & 71 & 2.32 {\footnotesize [1.69, 2.93]} & 33 & 77.1 & 27.3 \\

\bottomrule
\end{longtable}

\begin{longtable}{L{3.9cm}R{0.8cm}R{1.1cm}R{1.1cm}R{1.1cm}R{1.1cm}R{1.0cm}R{1.1cm}R{1.0cm}R{1.4cm}R{1.4cm}}
\caption{대분류별 점유율(\%), 상대 강도와 KL 대분류 간 항의 기여(지표 C)}\label{tab:ksco1}\\
\toprule
대분류 & 소분류 & 취업자 & C & $XS^{\mathrm{OE}}$ & $XS^{\beta}$ & $R^{\mathrm{C}}$ & $RX^{\mathrm{OE}}$ & $RX^{\beta}$ & KL 기여 (OE) & KL 기여 ($\beta$) \\
\midrule
\endfirsthead
% 04_16_analysis_kor_share_exposure_dist.do가 생성. 직접 고치지 말 것.
1 관리자 & 16 & 1.7 & 8.4 & 2.3 & 2.3 & 4.87 & 1.34 & 1.34 & -0.030 & -0.030 \\
2 전문가 및 관련 종사자 & 45 & 23.6 & 71.9 & 35.8 & 30.3 & 3.04 & 1.52 & 1.28 & -0.250 & -0.262 \\
3 사무 종사자 & 13 & 19.8 & 15.0 & 40.2 & 35.8 & 0.76 & 2.03 & 1.81 & +0.395 & +0.310 \\
4 서비스 종사자 & 13 & 11.4 & 0.6 & 1.9 & 6.4 & 0.05 & 0.16 & 0.56 & +0.021 & +0.152 \\
5 판매 종사자 & 5 & 10.8 & 1.2 & 16.6 & 10.3 & 0.11 & 1.53 & 0.95 & +0.434 & +0.222 \\
6 농림∙어업 숙련 종사자 & 4 & 5.9 & 0.0 & 0.0 & 2.5 & 0.00 & 0.00 & 0.42 & -0.000 & +0.162 \\
7 기능원 및 관련 기능 종사자 & 14 & 7.0 & 1.5 & 1.2 & 3.1 & 0.22 & 0.17 & 0.44 & -0.003 & +0.021 \\
8 장치∙기계 조작 및 조립 종사자 & 10 & 6.6 & 0.9 & 0.5 & 4.4 & 0.13 & 0.07 & 0.67 & -0.003 & +0.071 \\
9 단순노무 종사자 & 9 & 13.1 & 0.5 & 1.6 & 4.9 & 0.04 & 0.12 & 0.37 & +0.019 & +0.113 \\

\bottomrule
\end{longtable}
\end{small}

\paragraph{비유사성 지수.} 사용 분포를 노출도 기반 점유율과 같게 만들려면 사용 점유율의 60$\sim$63\%(OE)나 64$\sim$69\%(\,$\beta$)를 다른 직업으로 옮겨야 한다(표~\ref{tab:dist}). 같은 계산을 취업자 분포와 하면 사용은 73$\sim$77\%, 노출도 기반 점유율은 51\%(OE), 27\%($\beta$)다. 사용 분포가 ``취업자 비례''에서 벗어난 정도(73$\sim$77\%)를 노출도가 설명해 주는 부분은 크지 않다. 노출도를 고려해도 거리가 60$\sim$69\%로 줄 뿐이다. 이 거리의 67$\sim$72\%는 대분류 합으로만 계산해도 나온다. 즉 몫이 큰 이동은 대분류 사이에서 일어난다.

\paragraph{어느 대분류인가.} 표~\ref{tab:ksco1}에서 보이듯 전문가 대분류는 취업자의 20.8\%다. 노출도 기반 점유율은 35.4\%(OE)$\cdot$29.1\%($\beta$)이고, 사용 점유율은 71.9\%다. 사무 종사자는 노출도 기반 점유율이 39.7\%$\cdot$34.0\%로 가장 큰 대분류인데, 사용은 15.0\%다. 판매 종사자는 16.4\%$\cdot$9.8\% 대 1.2\%다. 상대 강도로 보면 사무는 노출도가 평균의 2.0$\sim$2.3배인데 사용은 0.87배이고, 판매는 노출도 1.04$\sim$1.73배에 사용 0.13배다. 반대로 관리자는 노출도 1.46$\sim$1.51배에 사용 5.57배다. paper 29의 ``전문가 +, 사무 $-$'' 패턴을 분포 거리로 다시 확인한 것이다. 이 노트는 거기에 크기를 더한다. KL 거리의 대분류 간 항 0.58(OE)$\cdot$0.76($\beta$)은 대부분 사무(+0.40$\cdot$+0.31)와 판매(+0.43$\cdot$+0.22)에서 나온다. $\beta$는 여기에 서비스$\cdot$농림$\cdot$단순노무(+0.11$\sim$+0.16)가 더해진다. KL은 사용 $> 0$인 129개 안에서 노출도 기반 점유율을 다시 정규화해 계산했다(사용 0인 소분류가 있으면 무한대가 되므로).

\paragraph{KL은 대분류 안이 크다.} KL 거리(1.35$\sim$2.32)의 대분류 간 몫은 31$\sim$35\%로 DI의 몫(67$\sim$72\%)보다 훨씬 작다. KL은 로그 비율에 무게를 둔다. 따라서 노출도 기반 점유율은 어느 정도 있는데 사용이 거의 없는 소분류에 크게 반응한다. 그런 소분류(451 조리사, 872 자동차 운전원, 311 기획 및 마케팅 사무원, 244 간호사 등, 3.5절)는 여러 대분류에 흩어져 있다. 두 거리를 함께 보면 이렇게 정리된다. 몫의 이동은 대분류 사이에서 일어나고, 비율의 어긋남은 대분류 안의 개별 직업에서 생긴다.

\subsection{탄력성: 노출도 차이가 사용에서 얼마나 증폭되는가}

\begin{small}
\begin{longtable}{L{2.4cm}L{0.8cm}R{2.1cm}R{2.1cm}R{0.9cm}R{0.8cm}R{1.5cm}R{1.1cm}R{1.1cm}R{1.1cm}}
\caption{$\ln R^m = \alpha + \beta \ln RX^x$의 기울기. 가중 = 취업자 가중 OLS(robust SE), 대괄호 = 부트스트랩 95\% 구간. 비가중 OLS와 분위수 회귀는 기울기만}\label{tab:elas}\\
\toprule
사용 지표 & 노출도 & 가중 OLS (SE) & 95\% 구간 & $R^2$ & $n$ & 비가중 OLS & $\tau$ = 0.1 & $\tau$ = 0.5 & $\tau$ = 0.9 \\
\midrule
\endfirsthead
% 04_16_analysis_kor_share_exposure_dist.do가 생성. 직접 고치지 말 것.
A 대화 & OE & 1.00 (0.23) & {\footnotesize [0.52, 1.52]} & 0.34 & 116 & 0.92 & 1.00 & 0.90 & 1.01 \\
 & $\beta$ & 1.80 (0.36) & {\footnotesize [1.26, 3.00]} & 0.30 & 129 & 2.31 & 1.88 & 2.45 & 2.27 \\
\midrule
B 절감 시간 & OE & 1.07 (0.26) & {\footnotesize [0.56, 1.66]} & 0.30 & 116 & 1.08 & 1.08 & 1.28 & 0.87 \\
 & $\beta$ & 2.21 (0.45) & {\footnotesize [1.54, 3.75]} & 0.32 & 129 & 2.84 & 2.38 & 3.12 & 2.72 \\
\midrule
C 취업자 환산 & OE & 1.06 (0.27) & {\footnotesize [0.55, 1.69]} & 0.29 & 116 & 1.08 & 1.10 & 1.12 & 0.94 \\
 & $\beta$ & 2.16 (0.44) & {\footnotesize [1.47, 3.66]} & 0.30 & 129 & 2.80 & 2.18 & 3.25 & 2.73 \\

\bottomrule
\end{longtable}
\end{small}

Observed exposure에 대한 탄력성은 1.00$\sim$1.07이다(구간 0.57$\sim$1.64, 표~\ref{tab:elas}). OE가 평균의 2배인 직업은 사용 강도도 평균의 약 2배라는 뜻이다. 사용 강도가 OE에 비례한다는 가설을 기각할 수 없다. 비가중(0.92$\sim$1.08)과 분위수 회귀(0.87$\sim$1.28)도 1 근처다. Eloundou $\beta$에 대한 탄력성은 1.80$\sim$2.21이고 구간(1.26$\sim$3.75)이 1을 포함하지 않는다. 비가중(2.31$\sim$2.84)과 분위수 회귀(1.88$\sim$3.25)에서는 더 크다. $\beta$가 평균의 2배인 직업은 사용 강도가 평균의 $2^{2.2} \approx 4.6$배다. 집중지수의 결론(OE 대비 차이 없음, $\beta$ 대비 더 쏠림)과 같다. Felten에 대한 준탄력성은 2.18$\sim$2.57이다(엑셀). Felten 1 단위(표준편차 약 0.8)가 오르면 사용 강도가 $e^{2.5} \approx 12$배가 된다. $R^2$은 0.29$\sim$0.41이다. 노출도가 사용 강도 로그 분산의 3분의 1 남짓만 설명한다. 탄력성은 사용 $> 0$인 129개(OE는 116개)로 추정하므로 점추정은 이전 판과 같고, 부트스트랩 구간만 163개 복원추출로 조금 바뀌었다.

분위수 회귀의 기울기는 분위마다 큰 차이가 없다($\beta$에 대해 0.1분위 1.88$\sim$2.38, 0.9분위 2.27$\sim$2.73). 증폭이 분포의 특정 구간에서만 일어나는 것이 아니라, 조건부 분포 전체가 함께 움직인다.

\subsection{직업별 거리}

\begin{small}
\begin{longtable}{L{6.6cm}R{1.6cm}R{1.6cm}R{1.6cm}R{1.6cm}R{1.8cm}}
\caption{지표 C와 노출도의 거리 분포. $d = \ln R^{\mathrm{C}} - \ln RX^x$, $\Delta z$ = 로버스트 $z$ 차이}\label{tab:dsum}\\
\toprule
항목 & $d$($\beta$) & $d$(OE) & $\Delta z$($\beta$) & $\Delta z$(OE) & $\Delta z$(Felten) \\
\midrule
\endfirsthead
% 04_16_analysis_kor_share_exposure_dist.do가 생성. 직접 고치지 말 것.
소분류 수 & 129 & 116 & 129 & 116 & 129 \\
10분위 & -4.30 & -3.47 & -0.87 & -1.05 & -1.10 \\
중위값 & -1.05 & -0.08 & 0.35 & 0.12 & 0.18 \\
90분위 & 2.25 & 2.29 & 1.48 & 1.04 & 0.89 \\
절댓값 평균 & 2.22 & 1.88 & 0.75 & 0.68 & 0.63 \\
\(|d| > \ln 2\)(2배 넘게 차이)인 소분류 수 & 109 & 93 & -- & -- & -- \\
\(|d| > \ln 10\)(10배 넘게 차이)인 소분류 수 & 51 & 31 & -- & -- & -- \\
순위 격차 RG와의 Spearman 상관 & 0.38 & 0.66 & 0.67 & 0.77 & 0.72 \\
취업자 수와의 Spearman 상관 & -0.57 & -0.64 & -0.38 & -0.49 & -0.50 \\
같은 노출도의 \(d\)와의 Spearman 상관 & -- & -- & 0.42 & 0.82 & -- \\

\bottomrule
\end{longtable}
\end{small}

\paragraph{거리의 크기.} 지표 C와 Eloundou $\beta$의 로그 거리는 절댓값 평균이 2.21이다(표~\ref{tab:dsum}). 전형적인 소분류에서 사용 강도와 노출도 강도가 $e^{2.2} \approx 9$배 차이 난다는 뜻이다. 사용 $> 0$인 129개 중 108개가 두 배 넘게, 51개가 열 배 넘게 차이 난다. 중위값은 $-1.00$이다. 중위 소분류의 사용 강도는 노출도가 예측하는 것의 약 3분의 1이다. 사용이 소수 직업에 몰려 있어 나머지 대부분이 평균 아래에 있기 때문이다. OE와의 거리는 조금 작다(절댓값 평균 1.88, 중위값 $-0.07$). 노출도 기반 점유율의 분모에 사용 0인 34개의 노출이 더해져 나머지 소분류의 $XS$가 작아졌으므로, $d$는 이전 판보다 $\beta$ 기준 약 0.05 커졌다.

\paragraph{순위 격차와의 관계.} $d$와 순위 격차 RG(paper 29 방법 1)의 순위상관은 $\beta$ 0.26, OE 0.61이다. 같은 직업이 순위로는 비슷해 보여도 거리로는 크게 다를 수 있다. 특히 $\beta$에서 그렇다. $\Delta z$는 퍼짐 차이를 지운 값이라 RG에 더 가깝다(0.70$\sim$0.79). $d$와 $\Delta z$의 순위상관은 $\beta$ 0.42, OE 0.82다. $\beta$에서 둘이 크게 다른 것은 $\beta$의 퍼짐(로그 분산 0.69)이 사용(10.6)보다 훨씬 작기 때문이다. 이 경우 $d$는 사실상 사용 강도 자체의 순서를 따른다. $d$와 취업자 수의 순위상관은 $-0.57$($\beta$), $-0.64$(OE)다. 큰 직업일수록 사용이 노출도에 못 미친다. 이는 사용 점유율이 직업 규모와 무관하다는 paper 32의 결과(규모 탄력성 0)와 같은 사실이다. 노출도 기반 점유율은 정의상 취업자에 비례하므로, 큰 직업에서 사용과의 거리가 벌어진다. RG는 163개 순위(사용 0은 평균 순위)로 계산하므로, 129개 순위로 계산하던 이전 판보다 $d$와의 상관이 낮다.

\begin{small}
\begin{longtable}{R{0.5cm}L{4.6cm}R{1.4cm}R{1.0cm}R{1.1cm}R{1.3cm}R{0.8cm}R{0.8cm}R{0.9cm}R{1.4cm}}
\caption{지표 C와 Eloundou $\beta$: 사용이 노출도보다 중심에서 더 위인 소분류($d$ 상위 15). 배 = $e^d$, RG = $r(\beta) - r(\mathrm{C})$}\label{tab:dpos}\\
\toprule
 & 소분류 & $R^{\mathrm{C}}$ & $RX^{\beta}$ & $d$ & 배 & $r$(C) & $r$($\beta$) & RG & 취업자(천 명) \\
\midrule
\endfirsthead
% 04_16_analysis_kor_share_exposure_dist.do가 생성. 직접 고치지 말 것.
1 & 213 생명 및 자연과학 관련 시험원 & 86 & 1.38 & +4.13 & 62.2 & 7 & 53 & +46 & 11 \\
2 & 149 기타 건설‧전기 및 생산 관련 관리자 & 69 & 1.19 & +4.06 & 58.0 & 34 & 69 & +35 & 3 \\
3 & 332 통계‧데이터 관련 사무원 & 74 & 2.45 & +3.41 & 30.2 & 6 & 6 & 0 & 16 \\
4 & 122 마케팅 및 광고‧홍보 관리자 & 46 & 1.55 & +3.40 & 29.9 & 18 & 29 & +11 & 9 \\
5 & 852 냉난방 관련 설비 조작원 & 10 & 0.51 & +2.99 & 19.9 & 35 & 121 & +86 & 17 \\
6 & 298 문화‧예술 관련 기획자 & 27 & 1.70 & +2.78 & 16.1 & 27 & 21 & -6 & 10 \\
7 & 159 기타 판매 및 고객 서비스 관리자 & 24 & 1.56 & +2.75 & 15.7 & 71 & 28 & -43 & 1 \\
8 & 291 작가 및 언론 관련 전문가 & 32 & 2.21 & +2.69 & 14.7 & 2 & 10 & +8 & 96 \\
9 & 292 학예사‧사서 및 기록물 관리사 & 21 & 1.50 & +2.65 & 14.2 & 12 & 37 & +25 & 31 \\
10 & 223 네트워크 및 정보 보안 전문가 & 29 & 2.18 & +2.60 & 13.5 & 8 & 11 & +3 & 33 \\
11 & 153 환경‧청소 및 경비 관련 관리자 & 19 & 1.53 & +2.53 & 12.5 & 51 & 33 & -18 & 4 \\
12 & 283 상품 기획‧홍보 및 조사 전문가 & 23 & 2.01 & +2.43 & 11.3 & 5 & 13 & +8 & 62 \\
13 & 269 기타 교육 전문가 & 14 & 1.43 & +2.31 & 10.0 & 13 & 47 & +34 & 46 \\
14 & 212 인문 및 사회과학 전문가 & 15 & 1.54 & +2.26 & 9.6 & 24 & 31 & +7 & 22 \\
15 & 282 금융 및 보험 전문가 & 14 & 1.66 & +2.11 & 8.2 & 14 & 23 & +9 & 47 \\

\bottomrule
\end{longtable}

\begin{longtable}{R{0.5cm}L{4.6cm}R{1.4cm}R{1.0cm}R{1.1cm}R{1.3cm}R{0.8cm}R{0.8cm}R{0.9cm}R{1.4cm}}
\caption{지표 C와 Eloundou $\beta$: 노출도가 사용보다 중심에서 더 위인 소분류($d$ 하위 15)}\label{tab:dneg}\\
\toprule
 & 소분류 & $R^{\mathrm{C}}$ & $RX^{\beta}$ & $d$ & 배 & $r$(C) & $r$($\beta$) & RG & 취업자(천 명) \\
\midrule
\endfirsthead
% 04_16_analysis_kor_share_exposure_dist.do가 생성. 직접 고치지 말 것.
1 & 451 조리사 & .00001 & 0.52 & -10.81 & 1/49423 & 129 & 106 & -23 & 813 \\
2 & 611 작물 재배 종사자 & .000085 & 0.41 & -8.48 & 1/4829 & 125 & 118 & -7 & 1,280 \\
3 & 872 자동차 운전원 & .00047 & 0.85 & -7.49 & 1/1784 & 119 & 77 & -42 & 1,012 \\
4 & 922 배달원 & .00055 & 0.73 & -7.19 & 1/1329 & 124 & 83 & -41 & 448 \\
5 & 630 어업 관련 종사자 & .0014 & 0.55 & -5.97 & 1/393 & 126 & 104 & -22 & 63 \\
6 & 952 음식 관련 단순 종사자 & .000064 & 0.02 & -5.47 & 1/237 & 128 & 129 & +1 & 398 \\
7 & 864 전기‧전자 부품 및 제품 조립원 & .0015 & 0.32 & -5.36 & 1/213 & 123 & 122 & -1 & 173 \\
8 & 873 물품 이동 장비 조작원 & .0027 & 0.50 & -5.22 & 1/186 & 122 & 109 & -13 & 139 \\
9 & 311 기획 및 마케팅 사무원 & .022 & 2.41 & -4.71 & 1/111 & 84 & 2 & -82 & 707 \\
10 & 244 간호사 & .0095 & 1.05 & -4.70 & 1/110 & 108 & 67 & -41 & 298 \\
11 & 620 임업 관련 종사자 & .0043 & 0.43 & -4.61 & 1/100 & 127 & 116 & -11 & 7 \\
12 & 613 축산 및 사육 관련 종사자 & .0073 & 0.55 & -4.31 & 1/75 & 117 & 105 & -12 & 95 \\
13 & 453 식음료 서비스 종사자 & .0061 & 0.45 & -4.30 & 1/74 & 106 & 114 & +8 & 514 \\
14 & 423 노인 및 장애인 돌봄 종사자 & .014 & 0.63 & -3.83 & 1/46 & 110 & 94 & -16 & 142 \\
15 & 751 자동차 정비원 & .006 & 0.24 & -3.71 & 1/41 & 115 & 126 & +11 & 164 \\

\bottomrule
\end{longtable}
\end{small}

\paragraph{사용이 노출도보다 위인 직업.} 213 생명 및 자연과학 관련 시험원은 사용 강도가 평균의 86배인데 노출도는 1.38배라, 사용이 노출도의 62배다(표~\ref{tab:dpos}). 149 기타 건설$\cdot$전기 및 생산 관련 관리자(58배), 332 통계$\cdot$데이터 관련 사무원(30배), 122 마케팅 및 광고$\cdot$홍보 관리자(30배)가 뒤를 잇는다. 15개 모두 취업자 10만 명 미만의 작은 직업이다. 순위로는 이 차이가 잘 보이지 않는다. 332는 사용 6위, $\beta$ 6위라 RG가 0이지만, 거리로는 3위다. 순위는 둘 다 상위라는 것만 말하지만, 거리는 사용이 노출도보다 한 자릿수 더 크다는 것까지 말한다. 반대로 852 냉난방 관련 설비 조작원은 RG가 +86으로 순위 격차가 가장 큰 축이지만, 거리로는 5위(20배)다. paper 29의 강건 목록과 겹치는 직업은 213, 852다.

\paragraph{노출도가 사용보다 위인 직업.} 아래쪽 끝은 사용 점유율이 거의 0인 직업이다(표~\ref{tab:dneg}). 451 조리사, 611 작물 재배, 872 자동차 운전원, 922 배달원, 630 어업, 952 음식 관련 단순 종사자 등이다. $\beta$는 이들 대부분에게 평균의 0.4$\sim$0.9배 노출을 주지만, 업무용 대화는 평균의 1천분의 1 이하다. 그중 451, 611, 620, 630, 864, 952는 AEI 공개 자료에 한 달만 나타나 값이 절반으로 계산된 소분류다(paper 25). 즉 $d$의 아래쪽 끝은 추정이 가장 불안정한 곳에 있다. 사용이 0인 34개는 이 목록에 들지 않는다($d$가 정의되지 않음). 이 점을 빼고 보면 눈에 띄는 직업은 311 기획 및 마케팅 사무원과 244 간호사다. 311은 $\beta$ 2위인데 사용은 84위다(RG $-82$, 사용이 노출도의 106분의 1). 244는 사용이 노출도의 105분의 1이다. 둘은 취업자가 30만$\sim$70만 명인 큰 직업이다.

\begin{footnotesize}
\begin{longtable}{R{0.4cm}L{5.4cm}L{5.4cm}L{5.4cm}}
\caption{지표 C의 세 거리별 상위$\cdot$하위 10 소분류(괄호 = 거리 값)}\label{tab:topC}\\
\toprule
 & $d$($\beta$) & $d$(OE) & $\Delta z$(Felten) \\
\midrule
\endfirsthead
% 04_16_analysis_kor_share_exposure_dist.do가 생성. 직접 고치지 말 것.
\multicolumn{4}{l}{\textit{사용이 노출도보다 중심에서 더 위(상위 10)}} \\
1 & 213 생명 및 자연과학 관련 시험원 (+4.13) & 149 기타 건설‧전기 및 생산 관련 관리자 (+7.10) & 213 생명 및 자연과학 관련 시험원 (+2.10) \\
2 & 149 기타 건설‧전기 및 생산 관련 관리자 (+4.06) & 213 생명 및 자연과학 관련 시험원 (+3.94) & 149 기타 건설‧전기 및 생산 관련 관리자 (+1.95) \\
3 & 332 통계‧데이터 관련 사무원 (+3.41) & 861 발전 및 배전 장치 조작원 (+3.32) & 852 냉난방 관련 설비 조작원 (+1.90) \\
4 & 122 마케팅 및 광고‧홍보 관리자 (+3.40) & 111 의회 의원‧고위 공무원 및 공공단체 임원 (+3.29) & 799 기타 기능 관련 종사자 (+1.72) \\
5 & 852 냉난방 관련 설비 조작원 (+2.99) & 237 환경공학‧가스‧에너지 기술자 및 시험원 (+3.02) & 730 목재‧가구‧악기 및 간판 관련 기능 종사자 (+1.31) \\
6 & 298 문화‧예술 관련 기획자 (+2.78) & 122 마케팅 및 광고‧홍보 관리자 (+2.95) & 292 학예사‧사서 및 기록물 관리사 (+1.20) \\
7 & 159 기타 판매 및 고객 서비스 관리자 (+2.75) & 332 통계‧데이터 관련 사무원 (+2.95) & 791 공예 및 귀금속 세공원 (+1.08) \\
8 & 291 작가 및 언론 관련 전문가 (+2.69) & 153 환경‧청소 및 경비 관련 관리자 (+2.82) & 226 통신 및 방송 송출 장비 기사 (+1.08) \\
9 & 292 학예사‧사서 및 기록물 관리사 (+2.65) & 226 통신 및 방송 송출 장비 기사 (+2.58) & 882 재활용 처리 및 소각로 조작원 (+0.99) \\
10 & 223 네트워크 및 정보 보안 전문가 (+2.60) & 159 기타 판매 및 고객 서비스 관리자 (+2.53) & 225 정보 시스템 및 웹 운영자 (+0.95) \\
\midrule
\multicolumn{4}{l}{\textit{노출도가 사용보다 중심에서 더 위(하위 10)}} \\
1 & 451 조리사 (-10.76) & 451 조리사 (-9.49) & 451 조리사 (-2.53) \\
2 & 611 작물 재배 종사자 (-8.43) & 922 배달원 (-6.48) & 531 통신 관련 판매 종사자 (-1.86) \\
3 & 872 자동차 운전원 (-7.44) & 872 자동차 운전원 (-4.94) & 952 음식 관련 단순 종사자 (-1.67) \\
4 & 922 배달원 (-7.14) & 864 전기‧전자 부품 및 제품 조립원 (-4.58) & 532 방문 및 노점 판매 종사자 (-1.54) \\
5 & 630 어업 관련 종사자 (-5.92) & 311 기획 및 마케팅 사무원 (-4.53) & 311 기획 및 마케팅 사무원 (-1.52) \\
6 & 952 음식 관련 단순 종사자 (-5.42) & 244 간호사 (-4.12) & 611 작물 재배 종사자 (-1.52) \\
7 & 864 전기‧전자 부품 및 제품 조립원 (-5.31) & 873 물품 이동 장비 조작원 (-3.69) & 245 영양사 (-1.28) \\
8 & 873 물품 이동 장비 조작원 (-5.17) & 331 회계 및 경리 사무원 (-3.58) & 872 자동차 운전원 (-1.22) \\
9 & 311 기획 및 마케팅 사무원 (-4.66) & 532 방문 및 노점 판매 종사자 (-3.57) & 322 운송 및 무역 사무원 (-1.20) \\
10 & 244 간호사 (-4.65) & 423 노인 및 장애인 돌봄 종사자 (-3.56) & 254 종교 관련 종사자 (-1.20) \\

\bottomrule
\end{longtable}
\end{footnotesize}

\paragraph{세 거리의 비교.} 표~\ref{tab:topC}에서 $d$(OE)의 위쪽에는 $d$($\beta$)에 없는 861 발전 및 배전 장치 조작원, 111 의회 의원$\cdot$고위 공무원, 237 환경공학$\cdot$가스$\cdot$에너지 기술자가 들어온다. OE가 낮은(0.01$\sim$0.03) 직업이라 작은 사용도 큰 비율이 된다. 아래쪽에는 331 회계 및 경리 사무원, 532 방문 및 노점 판매가 들어온다. OE가 높게 매긴 큰 사무$\cdot$판매 직업이다. $\Delta z$(Felten)은 퍼짐을 맞춘 비교라 극단값이 덜하다. 위쪽에 799 기타 기능, 730 목재$\cdot$가구, 791 공예 같은 기능직이 들어오는데, Felten이 이 직업들을 163개 중 116$\sim$119위로 낮게 매기기 때문이다. 아래쪽에는 531 통신 관련 판매, 245 영양사, 322 운송 및 무역 사무원이 있다. 세 거리 모두에 나오는 직업은 위쪽 213$\cdot$149, 아래쪽 451$\cdot$872$\cdot$311이다.

\section{순위 비교와 무엇이 다른가}

\begin{enumerate}[leftmargin=*]
  \item \textbf{크기.} 순위는 ``사용과 노출도의 순서가 0.73$\sim$0.75만큼 같다''(paper 28)까지만 말한다. 분포 비교는 사용의 직업 간 격차가 노출도보다 훨씬 크다는 것을 수치로 보여 준다(지니 0.90 대 0.37$\sim$0.64, 로그 분산 10.6 대 0.7$\sim$2.2). 일자리 영향의 직업 간 차이를 노출도로 예측하면 그 크기를 크게 과소평가하게 된다.
  \item \textbf{쏠림의 방향과 정도.} 사용이 고노출 직업에 몰린다는 것은 순위로도 알 수 있다. 그것이 노출도가 예측하는 정도인지, 그보다 더한지는 집중지수와 탄력성으로만 답할 수 있다. 답은 노출도 지표에 따라 다르다. Observed exposure에 대해서는 예측만큼이고(CI 차이 0, 탄력성 1), Eloundou $\beta$에 대해서는 두 배 넘게 증폭된다(탄력성 약 2). 다만 OE의 결과는 OE가 Anthropic 사용 자료를 일부 써서 만든 지표라는 점과 떼어 읽을 수 없다.
  \item \textbf{위치.} 같은 크기로 몰려 있어도(실효 직업 수 22 대 25) 몰린 곳은 다르다(DI 60\% 이상). 그 이동의 대부분은 사무$\cdot$판매에서 전문가$\cdot$관리자로의 대분류 간 이동이다.
  \item \textbf{어디서 생기는가.} Theil 분해로 보면, 노출도의 직업 간 차이는 주로 대분류 사이에서 생긴다(68$\sim$74\%). 사용의 차이는 주로 대분류 안에서 생긴다(57$\sim$60\%). 노출도가 대분류 안의 차이를 잘 구별하지 못한다는 뜻이다.
  \item \textbf{직업 목록.} 로그 거리 $d$는 순위 격차와 다른 직업을 고른다(순위상관 0.26$\sim$0.61). 순위 격차는 ``순서가 많이 바뀐 직업''(852, 311)을, 거리는 ``배율이 크게 다른 직업''(213, 149, 332, 451)을 고른다. 두 목록은 서로를 보완한다.
\end{enumerate}

\section{해석상 주의}

\begin{enumerate}[leftmargin=*]
  \item \textbf{아래 꼬리의 불안정.} 사용의 분산과 KL 거리, $d$의 하위 목록은 사용 점유율이 0에 가까운 소분류의 영향을 크게 받는다. 이 소분류들은 AEI 공개 하한 근처에 있거나 한 달만 공개된 값이다. 사용 0인 34개도 실제 사용이 0이라기보다 공개되지 않은 것이다. 지니$\cdot$DI는 이 영향이 작고, 로그 분산$\cdot$p90/p10$\cdot$KL$\cdot$$d$는 크다. 결론은 앞의 통계에 기대는 것이 안전하다.
  \item \textbf{연계표 균등 배분.} 소분류 점유율은 SOC 비중을 균등 가중치로 배분한 값이다(paper 21). 이 배분은 같은 SOC에 이어진 소분류들의 사용 강도를 기계적으로 다르게 만든다(취업자가 적은 쪽이 커짐). 그래서 대분류 안의 분산을 키우는 쪽으로 작용할 수 있다. 취업자 가중 배분 민감도 분석(paper 32의 과제)은 이 노트에서도 하지 않았다.
  \item \textbf{Observed exposure와의 순환성.} OE는 Eloundou 노출 등급과 AEI 사용 자료를 결합한 지표다\parencite{massenkoff-2026-labor-market-impacts-of-ai}. OE에 대한 ``CI 차이 0, 탄력성 1''은 일부 같은 자료에서 나온 결과다. 사용 자료와 독립적인 노출도로는 Eloundou $\beta$\parencite{eloundou-2023-gpts-are-gpts-an-early}와 Felten AIOE\parencite{felten-2021-occupational-industry-and-geographic-exposure,felten-2023-how-will-language-modelers-like-chatgpt}의 결과를 기준으로 삼는 것이 맞다.
  \item \textbf{부트스트랩 구간의 의미.} 구간은 소분류 표본 구성에 대한 민감도다. AEI 비중과 지역별고용조사의 표본 오차는 반영하지 않았다. 소분류는 모집단 전체(163개)이므로 이 구간을 통상의 표본 추론으로 읽으면 안 된다.
  \item \textbf{가중.} 주 결과는 취업자 가중이다. 비가중 지니는 사용 0.84$\sim$0.86, OE 0.64, $\beta$ 0.35이고 결론은 같다(엑셀 \texttt{disp} 시트). 비가중 탄력성은 가중보다 조금 크다.
  \item \textbf{노출도는 강도의 상한형 지표다.} $RX$는 ``과업 중 AI가 할 수 있는 비중''의 상대값이지 사용의 예측값이 아니다. 따라서 ``사용이 노출도보다 몰려 있다''는 노출도 지표가 틀렸다는 뜻이 아니다. 실제 채택이 노출된 과업 중 일부에서 훨씬 빠르게 일어난다는 뜻으로 읽어야 한다.
\end{enumerate}

\section{산출 파일}

\begin{itemize}[leftmargin=*]
  \item 코드: \texttt{code/04\_16\_analysis\_kor\_share\_exposure\_dist.do}(Stata 17). 입력은 \texttt{kor\_ai\_usage\_ksco3}(\texttt{03\_05\_merge\_kor\_ksco3\_ai\_usage.do}). 확인하는 항등식: 상대 강도의 가중 평균 1, $d = \ln(s / XS)$, CI($XS^x$) = 지니($RX^x$), Mata 가중 기울기 = Stata 가중 OLS. 지니$\cdot$집중지수$\cdot$Theil$\cdot$상위 몫은 Mata 함수로 계산하고, 부트스트랩에서도 같은 함수를 쓴다.
  \item 표: \texttt{results/table/33\_kor\_share\_exposure\_dist.xlsx}. 시트 \texttt{ksco3}(163개 소분류의 $w$, $s$, $XS$, $R$, $RX$, $d$, $\Delta z$, RG, $z$), \texttt{disp}(비가중 지니$\cdot$CV 포함), \texttt{hhi}, \texttt{conc}, \texttt{curve}, \texttt{dist}(반대 방향 KL 포함), \texttt{boot}(모든 통계의 점추정, 구간, 부트스트랩 평균$\cdot$SE), \texttt{elas}(Felten 준탄력성 포함), \texttt{dsum}(지표 A$\cdot$B 포함), \texttt{ksco1}.
  \item 이 노트의 표: \texttt{results/table/33\_disp.tex}, \texttt{33\_bdiff.tex}, \texttt{33\_conc.tex}, \texttt{33\_curve.tex}, \texttt{33\_hhi.tex}, \texttt{33\_dist.tex}, \texttt{33\_ksco1.tex}, \texttt{33\_elas.tex}, \texttt{33\_dsum.tex}, \texttt{33\_d\_pos.tex}, \texttt{33\_d\_neg.tex}, \texttt{33\_top\_C.tex}.
\end{itemize}

\printbibliography[title={참고문헌}]

\end{document}
